\renewcommand{\Anno}{1997}

\renewcommand{\NumeroProblema}{1}

\renewcommand{\Testo}{
Determinare tutte le coppie $(m,n)$ di numeri interi relativi tali che
$$m^3 -n^3 = 35.$$}

\renewcommand{\Risposta}{$(3,-2)$ e $(2,-3)$}

\renewcommand{\Aiutino}{
Scomporre $m^{3}-n^{3}$. Quali valori può assumere $m-n$?
}

\renewcommand{\SoluzioneUfficiale}{
Osserviamo innanzitutto che deve essere $m>n$.  Fattorizzando, possiamo riscrivere l'equazione nella forma
$$(m-n)(m^{2}+mn+n^{2})=35.$$

Poiché il termine $m-n$ è positivo, scomponendo 35 in fattori primi ci restano 4 possibilità, riassunte nella seguente tabella.
$$
\renewcommand{\arraystretch}{1.3}
\begin{array}{|c|c|}
	\hline
	m-n & m^{2}+mn+n^{2}  \\
	\hline\hline
	1 & 35  \\
	\hline
	5 & 7  \\
	\hline
	7 & 5  \\
	\hline
	35 & 1  \\
	\hline
\end{array}$$

Non ci resta quindi che distinguere 4 casi.
\begin{itemize}
	\item  Se $m=n+1$, sostituendo in $m^{2}+mn+n^{2}=35$ e svolgendo i 
	calcoli otteniamo l'equazione $3n^{2}+3n-34=0$, che non ha soluzioni 
	intere.

	\item  Se $m=n+5$, sostituendo in $m^{2}+mn+n^{2}=7$ e svolgendo i 
	calcoli otteniamo l'equazione $n^{2}+5n+6=0$, che ha come soluzioni 
	$n=-3$ (da cui $m=2$) ed $n=-2$ (da cui $m=3$).

	\item Se $m=n+7$, sostituendo in $m^{2}+mn+n^{2}=5$ e svolgendo i
	calcoli otteniamo l'equazione $3n^{2}+21n+46=0$, che non ha
	nemmeno soluzioni reali, avendo il discriminante negativo.

	\item  Se $m=n+35$, sostituendo in $m^{2}+mn+n^{2}=1$ e svolgendo i
	calcoli otteniamo l'equazione $3n^{2}+3\cdot 35n+35^{2}-1=0$, che non ha
	nemmeno soluzioni reali, in quanto il suo discriminante $9\cdot 
	35^{2}-12\cdot 35^{2}+12$ \`{e} chiaramente negativo.
\end{itemize}

Pertanto le uniche soluzioni intere sono le coppie $(3,-2)$ e $(2,-3)$.
}

\renewcommand{\SoluzioniAlternative}{
}

\renewcommand{\Osservazioni}{
}

\renewcommand{\Commenti}{
}